% =====================================================================
% Failure envelope. Phase 5 of RESEARCH_PLAN.md.
% Own section, not a limitations footnote.
% =====================================================================
\section{Failure Envelope: When DDI Loses}
\label{sec:failure-envelope}

Block-diagonal document masking makes a testable prediction of failure.
If documents cannot see each other, any question whose answer requires
composing information across two documents can be answered only by the
query and answer regions --- and only via whichever bridge the mask
permits. The honest version of this study measures that boundary and
reports it as the method's operating envelope.

\subsection{Single-hop versus multi-hop split}
\label{sec:failure-envelope-hop}

Every parity dataset admits a split into single-document-answerable
questions and questions that require cross-document composition. We
label each evaluation instance accordingly and report every mask cell
separately on the two splits. The single-hop split is where the parity
claim of Section~\ref{sec:results-quality} holds within $\Delta$ on
the primary instruct model. The multi-hop split is where DDI degrades,
and --- critically --- the degradation is not uniform across multi-hop
task classes. On the \texttt{ddi\_stress} families, \texttt{ddi\_qf}'s
F1 recovery (as a fraction of \texttt{full}'s F1) spans
$\approx 78\%$ on \texttt{needle} (single gold fact in a haystack of
distractors), $\approx 58\%$ on \texttt{contradiction},
$\approx 50\%$ on \texttt{coreference}, $\approx 28\%$ on
\texttt{fact\_chain}, and $\approx 12\%$ on \texttt{fact\_aggregation}
(combine facts spread across $N$ gold documents). Cross-document
aggregation is where DDI's mask most sharply blocks the required
mechanism; single-doc retrieval is where the mask barely registers.
The width of the operating envelope is therefore a task-family property,
not a single number. \todo{Family-level recovery is reported here at
one seed on the primary instruct model; a multi-seed re-run per family
would sharpen the numeric bounds, and a base-model rerun would test
whether the family ordering is instruct-specific.}

\begin{table}[t]
\centering
\caption{T5 --- quality on the single-hop versus multi-hop split of every
    parity dataset, per mask variant. The gap between the columns is the
    failure-envelope signal.}
\label{tab:t5}
\IfFileExists{generated/T5.tex}{%
    % Phase 3 document-count sweep. Source: results/phase3_n{2,5,10}.jsonl
\begin{tabular}{rccccc}
\toprule
& \multicolumn{2}{c}{EM} & & \multicolumn{2}{c}{F1} \\
$N$ & full & ddi\_qf & $\Delta$ & full & ddi\_qf \\
\midrule
2  & 0.62 & 0.55 & +0.07 & 0.80 & 0.71 \\
5  & 0.59 & 0.49 & +0.10 & 0.76 & 0.62 \\
10 & 0.46 & 0.46 & \textbf{0.00} & 0.64 & 0.61 \\
\bottomrule
\end{tabular}
%
}{%
    \todoreport{T5 hop-split table}{make split-hop}%
}
\end{table}

\figureifexists{generated/F4.pdf}{%
    F4 --- degradation vs.\ full baseline on the single-hop versus
    multi-hop split. A large gap between the two panels is the operating
    envelope of the method.%
}{fig:f4}

\subsection{Parallel Context Windows baseline}
\label{sec:failure-envelope-pcw}

Parallel Context Windows~\citep{ratner2023pcw} is the closest prior work
in the independent-context-window family and is the correct external
baseline for the parity claim. In DDI's mask family, PCW is exactly
\texttt{ddi\_qf} paired with the \texttt{restart} position layout; the
\texttt{phase5\_pcw} runner in the harness hard-codes that combination
so the PCW column of downstream tables is a first-class citizen rather
than a re-derivation of other rows. The two methods disagree on framing
--- PCW uses positional partitioning without an explicit mask family,
while DDI parameterizes the query--document bridge --- so a direct
comparison at matched shapes clarifies which effect is doing the work.
\todo{The PCW comparison rows have not yet been populated with
multi-seed measurements against DDI's \texttt{ddi\_qf} at matched
shapes on a shared parity dataset; the runner is in place but the
paired sweep has not been run.}

\subsection{Cross-document composition probe}
\label{sec:failure-envelope-composition}

The parity/multi-hop split relies on labels that inherit noise from the
dataset. As a controlled complement, we synthesize a cross-document
composition probe: pairs of documents whose combination is required for
the answer, with distractors and gold positions varied along the
Lost-in-the-Middle axes~\citep{liu2023lostmiddle}. Table~\ref{tab:t6} and
Figure~\ref{fig:f7} report quality as a function of the position and
count of the two gold documents; a monotone degradation in the
compositional axis is the signature we look for.

\begin{table}[t]
\centering
\caption{T6 --- synthetic composition probe. Quality as a function of the
    distance between the two gold documents and the number of distractors
    interposed.}
\label{tab:t6}
\IfFileExists{generated/T6.tex}{%
    \input{generated/T6.tex}%
}{%
    \todoreport{T6 composition table}{make composition-probe}%
}
\end{table}

\figureifexists{generated/F7.pdf}{%
    F7 --- composition-probe accuracy versus the distance between the two
    gold documents, one line per mask variant.%
}{fig:f7}
